%% APPENDIX E -- radio-frequency interference screening.
%% Owner: appendix-triage.  Carries \label{app:rfi}.
\section{Radio-frequency interference screening}
\label{app:rfi}

Interference is the first alternative explanation for an unattributed
crossing, so the screening this survey has, and the coverage it does not
have, are both set out here. Millimetre interferometry is an unusually poor
environment for terrestrial interference, though not an empty one: ITU Radio
Regulations Article~5 \citep{ITU2024} carries active allocations throughout
90--275\,GHz and WRC-19 identified parts of 275--450\,GHz for land-mobile and
fixed use. The specific hazard is the \RadarLo--\RadarHi\,GHz Earth
exploration-satellite allocation used by spaceborne cloud radar
\citep{CloudSat2002,EarthCARE2023}, a documented contaminant of ALMA Band~3;
\RadarNWin{} of the \RadarNTot{} searched windows overlaps it, carrying
\RadarNCross{} threshold crossings and \RadarNStage{} rank-passing events, so
no result here rests on it. What favours this band is that terrestrial
transmitters are sparse compared with centimetre wavelengths, that the array
sits at 5000\,m on a site chosen partly for radio quietness, and, decisively,
that an interferometer is not a total-power detector: a signal not in the far
field at the phase centre does not fringe-track but acquires a rapidly
varying geometric phase across the array, and is suppressed in the
correlation rather than imaged.

\emph{The coverage of the screening is uneven, and that is stated rather than
smoothed over.} The allocation test above and the sky- and
intermediate-frequency tests below are computed for every searched window
from the released catalogue, as is the \NCtrl-position rank screen, which
asks whether an excess is at the star or everywhere in the field. Recurrence
reaches only the crossings with a covering repeat block, and the visibility
fit is a consistency check rather than a filter. The correlator flagging is
the observatory's own, replayed from each block's delivered calibration
script, but it is not recorded per window in the release. \emph{No
per-channel interference flag of our own exists, and none can be recovered.}
Three measurements follow, and each is a test the population could have
failed.

\emph{First, clustering in sky frequency.} If interference dominated the
crossing population, the crossings would pile up at particular sky
frequencies. They do cluster --- but on the survey's own molecular lines.
The test needs a measured peak frequency in the frame in which an interfering
carrier would stand still, so it runs topocentrically over the
\AppMNCrossF{} of the \EvNCrossRaw{} crossings for which the release records
one, against the transitions the search was tuned to. For the other
\CsNCrossNoFreq, all \CsCrossNoFreqStar{} carbon-monoxide crossings, the
release gives the velocity offset from the line instead, and a frequency
reconstructed from a line offset cannot be used to ask whether frequencies
crowd onto lines. Of the \AppMNCrossF, \AppMTubeObs{} fall within
$\pm\DispoHalfKms$\,km\,s$^{-1}$ of a searched transition against
\AppMTubeNull{} expected if each window's peak fell anywhere on that
window's own channel grid, $\pv{\AppMTubeP}$. The densest band,
\AppMBandLo--\AppMBandHi\,GHz, holds \AppMBandN{} of them, \AppMBandNAttr{}
within that velocity window of CO($2{\to}1$); and once the line-attributed
crossings are set aside, the \AppMUnattrN{} that remain show no clustering at
all ($\pv{\AppMUnattrP}$). \emph{The concentration is the attribution, not a
residue left over after it.} Table~\ref{tab:ledger} counts the same band as
\EvTubeLedgerN{} crossings with \EvTubeLedgerAttr{} attributed, because it is
a different test over a different set --- all \EvNCrossRaw, each attributed
in its own star's rest frame --- and the \EvTubeExtra{} it adds are exactly
the ones whose peak frequency this test has no measurement of.

\emph{Second, fixed intermediate frequency.} An artefact of the signal chain
would sit at a fixed position within the spectral window rather than at a
fixed sky frequency. This test is run over every window for which the
catalogue releases a peak frequency --- \IfPopNWin{} of the \NWindows{}, of
which only \IfPopNCross{} carry a threshold crossing --- because such an artefact would
occupy the same intermediate frequency in all of them and not only where the
trigger fires. It does not: the peaks' fractional positions within their own
windows have median \RfiFracMed{} and are uniform.

\emph{Third, a cross-target occupancy screen over every searched window and
not only the crossings.} Interference belongs to the observatory and not to
the star, so a terrestrial carrier should recur at the same sky frequency
toward \emph{unrelated} stars. It does not even for the leading event:
\RfiEventLead; \RfiOtherWin{} further windows, toward \RfiOtherStarsRange{}
other stars in \RfiOtherBlocksRange{} other blocks, cover \RfiFreqThem, and
\RfiSameChanClause.\RfiNotTestableClause{} Over all \RfiOccNWin{} searched
windows toward \RfiOccNStars{} stars, with each window's peak randomised on
its own channel grid, the number of cross-star coincidences within
\RfiOccTol\,MHz is \RfiOccObs{} against
\RfiOccNull$\,\pm\,$\RfiOccNullSd{} expected ($\pv{\RfiOccP}$); among the
\RfiOccSigN{} windows whose peak exceeds the trigger it is \RfiOccSigObs{}
against \RfiOccSigNull$\,\pm\,$\RfiOccSigNullSd, i.e.\ \emph{below} the
null, and with the molecular-line mask applied \RfiOccMskObs{} against
\RfiOccMskNull$\,\pm\,$\RfiOccMskNullSd. There is no excess to explain. Two
features of the screen matter. It is keyed on the star and not on the product
directory, because a star's own repeat observations are not different
targets; and \RfiOccMaskN{} of those \RfiOccSigN{} windows sit on a masked
transition (\RfiOccSpeciesList), so the mask has to be applied first --- a
carbon monoxide line falls at nearly the same sky frequency toward every
nearby star, which is exactly the coincidence the screen is looking for.

\emph{One execution block fails these tests, and it fails them as a block
rather than as a carrier.} Its \DqBlockNCross{} crossings toward
\LgVerRejectStar{} are the highest statistics in the survey outside
$\beta$~Pictoris and HD~48370, and the field statistics that identify the
failure are in \S\ref{sec:bdblock}. What belongs here is whether
interference accounts for it, and three readings say not. The visibility fit
returns a large imaginary part, placing the emission away from the phase
centre. The block's \BdNSibling{} sibling executions of the same field, with
the same array and the same spectral setup, produce no crossing at all. And
it is not a carrier at a fixed sky frequency: each of the four frequencies is
covered by \DqOtherWinLo--\DqOtherWinHi{} further windows toward
\DqOtherStarLo--\DqOtherStarHi{} other stars, and \DqOtherSame{} of them
carries a crossing in the same channel.

\emph{One reading does point at a mechanism, and it can be tested across the
survey.} \DqBlockIfMult{} of the \DqBlockNCross{} sit at the identical
channel of two different spectral windows, which is the signature of a fixed
intermediate frequency --- and a feature with no geometric fringe phase, a
spur common to every antenna being the obvious case, does not fringe-track,
so it gathers at the phase centre rather than at the star. Here the two are a
twentieth of an arcsecond apart, so the block cannot test its own hypothesis;
the survey can, because the star lies \EvPcOffLo--\EvPcOffHi\,arcsec from the
phase centre in \EvPcNWin{} windows whose control annulus brackets it,
putting a median of \EvPcNNear{} probes within two synthesised beams of the
phase centre and the star well outside. Against the probes at the same
distance from the star but elsewhere in azimuth --- matching the radius
matters, or a centrally concentrated disc would imitate the effect --- those
probes are higher by $\EvPcDiff\pm\EvPcErr$ in the statistic, consistent with
zero and bounded by \EvPcBound, smaller by a factor of about \EvPcFactor{}
than the \EvPcNeed{} by which this block's ring is raised. \emph{The phase
centre is not where the excess lives.}
\textbf{The evidence supports a
data-quality failure of one execution block, and that is how these
\EvNCrossFlag{} crossings are dispositioned} (\S\ref{sec:bdblock}): the
criterion of Appendix~\ref{sec:blockquality} flags \DqNWin{} of the
\NWindows{} windows, toward \DqNStar{} stars (\DqStars), the second being a
resolved disc rather than a defect.

\emph{None of this is a substitute for the ON/OFF cadence a dedicated search
uses, and we do not present it as one.} We claim only that \emph{no
identified interference mechanism accounts for the unattributed events}:
allocation tables cannot exclude harmonics, intermodulation products,
unconventional emitters or near-field sources, and archival data without a
simultaneous off-source strategy cannot do better. The three tests establish
that the crossing population as a whole carries none of the signatures
interference leaves, which is a statement about the population and not a
clearance of any individual crossing.

%% ------------------------------------------------------------------
%% NOTE: per 06_discussion.tex's own R6, this appendix no longer ends on
%% "the null result resting on internal validation alone" -- that statement
%% now lives once, in the discussion's lessons subsection.
%% REQUEST (appendix-triage -> integrator):
%%   fig:appm (the power comparison of the per-channel multiplicity statistic
%%   against the scan statistic) is deleted: it exists to compare our own two
%%   statistics, and the paper now states only the one it uses.  The
%%   \AppMPowerOld* and \RfiOcc{Prior,Dir,Stale}* macro families are no
%%   longer referenced and will be retired.  figures/appm_power.pdf is no
%%   longer included.
%% ------------------------------------------------------------------
